% =====================================================================
%  Integumentary Emergence in the Jammed-Granular-Vacuum Framework
%  VP Theory -- single-author whitepaper  (paper_id: integumentary_vp_site)
%  Author: Young Jae Lee (ORCID 0009-0002-7535-8245)
%  Version 1.0.0  -- build 2026-06-19
%  Concept DOI: 10.5281/zenodo.20754541 (all versions) | version DOI: 10.5281/zenodo.20754542 (v1.0.0)
%  License: CC BY 4.0
%  Compile: pdflatex (run twice for the table of contents / references)
% =====================================================================
\documentclass[11pt,a4paper]{article}

\usepackage[utf8]{inputenc}
\usepackage[T1]{fontenc}
\usepackage{mathptmx}
\usepackage[scaled=0.92]{helvet}
\usepackage[margin=2.4cm]{geometry}
\usepackage{amsmath,amssymb}
\usepackage{booktabs}
\usepackage{array}
\usepackage{longtable}
\usepackage{textcomp}
\usepackage{microtype}
\usepackage{xcolor}
\usepackage{enumitem}
\usepackage{titlesec}
\usepackage{caption}
\usepackage[hidelinks,colorlinks=true,linkcolor=linknavy,urlcolor=linknavy,citecolor=linknavy]{hyperref}

% ---- colours -------------------------------------------------------
\definecolor{linknavy}{HTML}{1A4E8A}
\definecolor{gV}{HTML}{1F7A3D}
\definecolor{gL}{HTML}{1A5FA8}
\definecolor{gF}{HTML}{8A5A00}
\definecolor{gO}{HTML}{9A2A2A}
\definecolor{rule}{HTML}{B9C2CC}
\definecolor{boxbg}{HTML}{F4F6F8}

% ---- grade badges --------------------------------------------------
\newcommand{\V}{\textcolor{gV}{\textbf{[V]}}}
\newcommand{\Lg}{\textcolor{gL}{\textbf{[L]}}}
\newcommand{\Fg}{\textcolor{gF}{\textbf{[F]}}}
\newcommand{\Og}{\textcolor{gO}{\textbf{[O]}}}

% ---- glyph shorthands (ASCII-safe; pdflatex) -----------------------
\newcommand{\gam}{\ensuremath{\gamma}}
\newcommand{\ra}{\ensuremath{\rightarrow}}
\newcommand{\mul}{\ensuremath{\times}}
\newcommand{\apx}{\ensuremath{\approx}}
\newcommand{\cd}{\ensuremath{\cdot}}
\newcommand{\dg}{\ensuremath{^{\circ}}}
\newcommand{\um}{\ensuremath{\mu}m}
\newcommand{\tewlunit}{\ensuremath{\mathrm{g\,m^{-2}\,h^{-1}}}}

% ---- section styling ----------------------------------------------
\titleformat{\section}{\large\bfseries\sffamily}{\thesection.}{0.6em}{}
\titleformat{\subsection}{\normalsize\bfseries\sffamily}{\thesubsection}{0.6em}{}
\titlespacing*{\section}{0pt}{1.6ex plus .3ex}{0.8ex}
\setlength{\parskip}{0.45em}
\setlength{\parindent}{0pt}
\renewcommand{\arraystretch}{1.18}
\captionsetup{font=small,labelfont=bf}

% answer-first callout box
\newcommand{\answer}[1]{%
  \par\vspace{0.2em}\noindent
  \fcolorbox{rule}{boxbg}{\parbox{\dimexpr\linewidth-2\fboxsep-2\fboxrule}{%
  \small\textit{#1}}}\par\vspace{0.4em}}

\hypersetup{
  pdftitle={Integumentary Emergence in the Jammed-Granular-Vacuum Framework},
  pdfauthor={Young Jae Lee},
  pdfsubject={VP Theory -- physics-derived emergence of the integumentary system on one bistable substrate},
  pdfkeywords={VP Theory, jammed granular vacuum, integumentary system, epidermal barrier, wound healing, jamming unjamming, melanin photoprotection, UV carcinogenesis, hysteresis, reproducibility}
}

\begin{document}

% =====================================================================
%  TITLE
% =====================================================================
\begin{titlepage}
\centering
{\sffamily\bfseries\LARGE
Integumentary Emergence in the\\[2pt]
Jammed-Granular-Vacuum Framework\par}
\vspace{0.7em}
{\sffamily\large
The Epidermal Barrier, Wound-Healing Unjamming,\\
and the Ultraviolet-Carcinogenesis Showcase\par}
\vspace{1.4em}
{\large Young Jae Lee\par}
\vspace{0.2em}
{\small Independent Researcher \;\textbar\; ORCID \href{https://orcid.org/0009-0002-7535-8245}{0009-0002-7535-8245}\par}
\vspace{0.2em}
{\small VP Theory --- Jammed-Granular-Vacuum Emergence \;\textbar\; \href{https://jamming-physics.org/integumentary/}{jamming-physics.org/integumentary}\par}

\vspace{1.8em}
\begin{minipage}{0.86\textwidth}\small\centering
\renewcommand{\arraystretch}{1.25}
\begin{tabular}{r l}
\textbf{Package} & \texttt{integumentary\_vp\_site} \;(code: \texttt{skn})\\
\textbf{Physical class} & jamming --- barrier/interface + external insult\\
\textbf{Version} & 1.0.0 \;(in-lane program complete) --- build 2026-06-19\\
\textbf{Governance} & VP-SPEC v1.8 (constitution C0 overrides on conflict)\\
\textbf{DOI (concept)} & \href{https://doi.org/10.5281/zenodo.20754541}{10.5281/zenodo.20754541} \;(all versions)\\
\textbf{DOI (this version)} & \href{https://doi.org/10.5281/zenodo.20754542}{10.5281/zenodo.20754542} \;(v1.0.0)\\
\textbf{License} & Creative Commons Attribution 4.0 (CC BY 4.0)\\
\end{tabular}
\end{minipage}

\vspace{1.6em}
\begin{minipage}{0.9\textwidth}\small
\textbf{Abstract.} The integumentary system --- the body's boundary --- is reconstructed not as a designed
organ set but as an \emph{emergent} one. In the VP (jammed-granular-vacuum) framework every cell programme rides one
vendored bistable primitive, the R19 switch $\dot s = g\,s - s^{3} + h$, whose control parameter $g$ is a
\emph{measured} master-gene order parameter $\gamma$ carried read-only from the DNA gene-clock. Four organs
(epidermis TP63, keratinocyte KRT14, melanocyte MITF, skin-appendage EDAR) switch on from their measured $\gamma$,
and a fifth (sebaceous gland PRDM1) is added by the same promoter-thermodynamics pipeline that reproduces MITF
and EDAR byte-for-byte; the developmental order is read directly off $\gamma$, never fitted. On this one substrate
the work verifies ten discriminant targets --- the epidermal water barrier, wound healing as an unjamming
transition around the vertex-model shape index $q^{*}=3.81$, melanin as an ultraviolet negative feedback, epidermal
turnover as a dwell-cascade, sweating as a regulated interface flux, ultraviolet carcinogenesis as a convex
multistage rate, and four additive hysteretic dynamics (the hair-follicle relaxation oscillator and the sebaceous,
adhesion, and vasomotor jams) --- and reproduces twenty-three skin diseases as \emph{signed perturbations of one
existing knob}, each established treatment being the same knob reversed. The decisive internal test is that
clinically opposite entities (psoriasis/ichthyosis, melasma/vitiligo, pemphigus/pemphigoid, rosacea/Raynaud)
fall out of the same switch driven in opposite directions, with no constant changed between them. The programme
holds a strict no-tuning rule --- every constant is a measured input or substrate-derived, never chosen to hit a
target --- and is reproducible to seven byte-identical SHA-256 hashes that a single release-audit script re-verifies
($7/7$ match, all gates green). Claims are graded honestly throughout: \V\ simulation-verified shape, \Lg\ cited
literature anchor, \Fg\ substrate-forced regime scale, \Og\ absolute magnitude open with a stated obstacle.
\end{minipage}

\vfill
{\footnotesize This whitepaper is the human-readable companion to the canonical per-title HTML at
\href{https://jamming-physics.org/integumentary/}{jamming-physics.org/integumentary}. The reproducible engine,
governance documents, and this source are bundled in the accompanying reproducibility archive.\par}
\end{titlepage}

\tableofcontents
\vspace{1em}
\hrule
\vspace{0.6em}

% =====================================================================
\section*{Conventions and the grading system}
\addcontentsline{toc}{section}{Conventions and the grading system}

Every result in this document is reproduced from one vendored primitive --- the R19 bistable switch
$\dot s = g\,s - s^{3} + h$, with its spinodal, basin barrier, and dwell time --- driven by the measured
master-gene order parameters $\gamma$. The governing discipline is the \textbf{no-tuning invariant}:
\emph{every constant is either a measured input (cited and locked) or substrate-derived; none is chosen to fit a
target}. The methodology is \textbf{LOCK $\to$ Derive $\to$ Gate}: inputs are locked, consequences are derived, and a
gate suite refuses to let the writing phase proceed until the research battery is green. The pipeline is
deterministic --- BLAS pinned single-thread, all grids fixed, no random-number generator --- so two independent runs
produce a byte-identical result hashed to SHA-256.

Claims carry one of four grades, used uniformly:

\begin{description}[leftmargin=2.6em,style=nextline,itemsep=1pt]
\item[\V\ Simulation-verified] a shape, control, ordering, sign, or discontinuity reproduced by the substrate.
\item[\Lg\ Cited anchor] a measured number taken from the external literature.
\item[\Fg\ Forced (substrate)] a dimensionless regime-scale set-point forced by the substrate (e.g.\ a spinodal
      computed from measured $\gamma$); forced, not fitted.
\item[\Og\ Open (obstacle stated)] an absolute magnitude the package does \emph{not} yet derive, each listed with
      the specific external input that would close it. An \Og\ without a stated obstacle is a gate failure.
\end{description}

% =====================================================================
\section{Integumentary emergence: the body boundary from measured $\gamma$}

\answer{Four skin organs emerge from measured master-gene $\gamma$ on the shared R19 bistable substrate. Ordered
by ascending $\gamma$ the predicted developmental sequence is epidermis $\to$ skin-appendage $\to$ melanocyte
$\to$ keratinocyte --- a consequence of the measured inputs, not a fit.}

The integumentary system is the organism's interface with the world: a barrier that holds water in, a screen that
holds ultraviolet out, a sensor sheet, and a thermostat. Here it is not designed; it \emph{emerges} when a small
set of master genes switch on the shared bistable substrate. Each organ rides the same vendored switch, whose
control parameter $g$ is the measured master-gene $\gamma$ taken read-only from the DNA gene-clock package.
Identity and developmental order follow from those measured numbers; nothing is tuned to a skin phenotype.

\begin{table}[h]
\centering\small
\caption{The integumentary organ atlas. $\gamma$ is a measured, read-only input; the spinodal is substrate-derived
from $\gamma$ and sets the threshold at which each state flips discontinuously. PRDM1 was fetched, cached, and
vendored in v0.6.0 through the identical promoter-thermodynamics pipeline validated against MITF and EDAR, and
reproduces offline bit-for-bit ($\gamma=1.343156$).}
\begin{tabular}{l l c c l}
\toprule
\textbf{Organ} & \textbf{Master gene} & \textbf{$\gamma$ (measured)} & \textbf{Spinodal} & \textbf{Role class}\\
\midrule
Epidermis        & TP63          & 1.3643 & 0.6134 & barrier\\
Keratinocyte     & KRT14         & 1.4894 & 0.6996 & barrier\\
Melanocyte       & MITF          & 1.3945 & 0.6338 & defense\\
Skin appendage   & EDAR          & 1.3696 & 0.6169 & appendage\\
Sebaceous gland  & PRDM1 (Blimp1)& 1.3432 & 0.5990 & jamming\\
\bottomrule
\end{tabular}
\end{table}

Ordering the four core organs by ascending $\gamma$ gives the substrate's prediction for the sequence in which
they consolidate: \textbf{epidermis $\to$ skin-appendage $\to$ melanocyte $\to$ keratinocyte}. The epidermal field
commits first, appendages and melanocytes follow, and the stiff keratinocyte programme last. This ordering is a
falsifiable consequence of the measured $\gamma$, not an input \V. PRDM1, committed on its biology \emph{before}
$\gamma$ was read, then turns out to carry the \emph{lowest} $\gamma$ in the whole atlas, so its functional spinodal
opens earliest --- an emergence-order prediction graded by sign \V.

% =====================================================================
\section{Epidermal barrier and transepidermal water loss (T1)}

\answer{Transepidermal water loss rises as $1/N$ with barrier thickness, crossing twice baseline at half thickness,
then collapses discontinuously when a chemical insult drives the keratinocyte switch past its spinodal. Insult
tolerance scales as $\gamma^{1.5}$, ranking the keratinocyte highest.}

The stratum corneum is a diffusion barrier and transepidermal water loss (TEWL) is the steady flux across it. For
a barrier of $N$ layers the flux scales as $1/N$, so thinning raises TEWL smoothly and monotonically; on a 15-layer
reference barrier the flux crosses twice baseline once thickness halves --- here at 7 layers. That is the ordinary
graded failure mode. \emph{Insult} is different. A chemical insult does not thin the barrier; it drives the
keratinocyte switch field $h$ toward its spinodal. The ON-basin depth is robust until the drive reaches a fraction
$\apx 1.0033$ of the spinodal, and then the barrier collapses \emph{discontinuously}, TEWL jumping rather than
drifting \V. The absolute insult a layer absorbs before that snap scales as $\gamma^{1.5}$, so the organs rank in
resistance as $\mathrm{KRT14}>\mathrm{MITF}>\mathrm{EDAR}>\mathrm{TP63}$ \V; the keratinocyte absolute tolerance is
0.70 dimensionless drive units. The conversion to a clinical TEWL in \tewlunit\ is open \Og: it needs the
stratum-corneum lipid permeability and the trans-barrier water-activity gradient as cited inputs.

% =====================================================================
\section{Wound healing as jamming, unjamming, and re-jamming (T2)}

\answer{Re-epithelialisation is an unjamming--rejamming cycle around the vertex-model shape index $q^{*}=3.81$. A
free edge unjams ($q$ rises to 4.07), cells migrate, confluence re-jams the sheet ($q$ falls to 3.55); below a
critical drive the wound fails to close.}

A confluent epithelial sheet is a \emph{jammed} solid: cells are caged by neighbours and cannot rearrange. Its
jammed-versus-fluid state is set by a single geometric control, the cell shape index $q$, with a vertex-model
threshold at $q^{*}=3.81$ (Bi, Lopez, Schwarz \& Manning, 2015--2016) \Lg: below it the sheet is jammed, above it
it flows. Injury frees an edge; the free-edge cue drives edge cells across $q^{*}$, \emph{unjamming} them (here
$q_{\max}=4.07$), and the unjammed cells migrate collectively into the gap. As the gap closes the free-edge cue
vanishes, confluent crowding pushes the edge field back below the spinodal, and the sheet \emph{re-jams}
($q_{\text{final}}=3.55$). The same model predicts failure: below a critical unjamming drive of about 0.15 the edge
never crosses $q^{*}$, the sheet stays jammed, and the gap does not close --- the substrate's picture of a chronic,
non-healing wound \V. The absolute closure rate in \ensuremath{\mathrm{\mu m\,h^{-1}}} is open \Og: it needs a
measured single-cell migration speed.

% =====================================================================
\section{Melanin photoprotection as ultraviolet negative feedback (T3)}

\answer{Ultraviolet drives MITF$\to$melanin, and melanin screens ultraviolet, closing a negative-feedback loop.
The melanin response is concave and plateauing; delivered ultraviolet is attenuated to about 21\%, and removing the
feedback abolishes the protection.}

Ultraviolet light drives the melanocyte master gene MITF, which makes melanin; melanin then absorbs ultraviolet ---
including the ultraviolet that triggered it. Synthesis and screening therefore form a negative-feedback loop,
$h_{\text{eff}} = \mathrm{UV}\cdot e^{-kM}$. The self-consistent steady melanin level is consequently
\emph{concave} in dose, rising toward a plateau rather than tracking the dose linearly: across the swept range the
high-dose level is only about $1.09\mul$ the mid-dose level, the signature of a regulated quantity. The ultraviolet
that actually reaches the DNA below is the screened flux $\mathrm{UV}\cdot e^{-kM}$, attenuated to about 21\% of the
incident dose at the top of the range --- substantial but partial photoprotection (a tan is not sunscreen). The
control experiment removes the feedback ($k\to0$): delivered-ultraviolet attenuation returns to $1\mul$ (no
protection at all), identifying the melanin screen --- not any synthesis cap --- as the cause \V. The absolute
minimal erythema dose and melanin optical density are open \Og: they need the melanin extinction coefficient.

% =====================================================================
\section{Epidermal turnover as a substrate dwell-cascade (T4)}

\answer{Total turnover is the sum of compartment residences. Viable epidermis and stratum corneum share TP63, so
the substrate gives them equal dwell; with the cited $\apx 14$-day corneum transit the total is $\apx 28$ days,
inside the cited 28--40-day window.}

Epidermal renewal moves a cell from the basal layer to the surface and off; the total turnover time is therefore a
\emph{sum of compartment residence times}, long established by labelled-cell studies. The two compartments here ---
viable epidermis and stratum corneum --- share the same master gene TP63, so the substrate assigns them the same
dwell time (the dwell ratio is 1): equal $\gamma$ gives equal substrate dwell, a prediction not a fit. With the
stratum-corneum transit fixed to its cited 14 days \Lg, the equal-dwell viable phase is also 14 days and the total
epidermal turnover is 28 days, inside the cited 28--40-day window for young-adult skin \V. Because dwell grows with
$\gamma$, the four organs rank in dwell as $\mathrm{KRT14}>\mathrm{MITF}>\mathrm{EDAR}>\mathrm{TP63}$. The absolute
basal cycle time per cell is open \Og: it needs a per-cell calibration. (See \S\ref{sec:anchors}: current literature
puts whole-epidermis turnover at $\apx 28$--48 days, extending to $\apx 56$ days in newer estimates, so the
modelled 28 days sits at the conservative lower edge of the window --- stated, not tuned.)

% =====================================================================
\section{Thermoregulation: the sweat gland as a regulated interface flux (T5)}

\answer{A thermal load raises core temperature until it crosses the EDAR sweat-switch threshold; the evaporative
interface flux then regulates temperature, dropping the temperature--load slope from 2.00 to 0.77, until capacity
saturates and temperature runs away.}

A thermal load raises core temperature. The sweat gland is governed by the EDAR switch and stays off until the
thermal drive clears its threshold (recruitment begins at a load of about 0.3); below that the skin sheds heat only
passively. Once recruited, the gland secretes in proportion to the temperature error and evaporative cooling
carries heat away as an interface flux --- a feedback controller engaging at a set-point, not a linear radiator. The
regulation is visible in the slope of core temperature against load: about 2.00 before sweating, falling sharply to
about 0.77 once the evaporative flux engages, the controller actively flattening the rise. Above the maximum sweat
capacity the flux saturates, the slope climbs back to about 2.00, and temperature runs away --- the substrate's
picture of the heat-stroke limit \V. The absolute set-point (37\dg C) and the absolute sweat rate are open \Og:
they need per-gland output and body heat capacity.

% =====================================================================
\section{Ultraviolet carcinogenesis: the squamous/melanoma dichotomy}

\answer{A single convex, multistage malignant rate splits the two ultraviolet cancers. Squamous-cell carcinoma
rises near-linearly with cumulative dose ($r=0.999$); melanoma is intermittent-sensitive (burst relative risk up to
$\apx 5.7\mul$); and a chronic-exposure tan protects against melanoma ($\apx 2.6\mul$).}

In this framework a carcinogen is a sustained aberrant drive that lowers the barrier between the normal and
malignant basins of the same cell switch; the malignant-crossing rate is Kramers-like,
$\text{rate}\sim e^{|h|/\text{spinodal}}$. Because malignancy needs several independent hits, the realised rate is
multistage, $\text{rate}^{K}$ with $K\apx 5$ (Armitage--Doll) \Lg --- a cited biological structure, not a fitted
knob. That single convex rate splits the two ultraviolet cancers.

\textbf{Squamous-cell carcinoma} tracks cumulative dose. At a fixed chronic intensity the cumulative
malignant-initiation probability accumulates linearly in total dose, so SCC rises near-linearly with cumulative
ultraviolet (linear to $r=0.999$) \V, matching the recognised near-linear dependence of cutaneous SCC on
cumulative, occupational ultraviolet.

\textbf{Melanoma} tracks intermittent bursts. For a fixed cumulative dose, concentrating it into bursts raises the
integrated hazard, because a convex rate rewards peaks (Jensen's inequality). Delivering the same dose
intermittently raises the melanoma relative risk by up to about $5.7\mul$ here \V --- the substrate's account of
melanoma's sunburn/intermittent-exposure sensitivity.

\textbf{The chronic-exposure paradox.} A spread, chronic dose builds a melanin screen (a tan) that a fast sunburn
cannot, lowering the effective intensity; here the tan reduces the chronic-arm hazard by about $2.6\mul$ \V,
reproducing the otherwise puzzling inverse association between chronic occupational ultraviolet and melanoma. The
cited anchor is Gandini et al.\ (2005), a meta-analysis of 57 studies: intermittent exposure raises melanoma risk
(summary relative risk 1.61, 95\% CI 1.31--1.99) while chronic/occupational exposure does not, and sunburn roughly
doubles risk \Lg. Absolute incidences and relative-risk magnitudes are open \Og: they need a population baseline
rate and an absolute dose calibration.

% =====================================================================
\section{Methods, grades, and reproducibility}\label{sec:methods}

\answer{All results derive from the R19 bistable substrate and measured master-gene $\gamma$ under a strict
no-tuning rule. The pipeline is bit-reproducible to an identical SHA-256, and every claim is graded, with each open
item carrying a stated obstacle.}

Every result is reproduced from one vendored primitive, the R19 bistable switch $\dot s = g\,s - s^{3} + h$, with
its spinodal, basin barrier, and dwell. The only organism-specific inputs are the measured master-gene $\gamma$
values carried read-only from the DNA gene-clock. The governing rule is the no-tuning invariant: every constant is
a measured input (cited and locked) or substrate-derived, never chosen to fit a target. The measured $\gamma$ drive
falsifiable cross-checks --- the developmental order, the $\gamma^{1.5}$ insult ordering, the equal-dwell turnover
--- while the dimensionless dynamical scales set the regime in which each mechanism is exhibited and are graded
honestly rather than calibrated to data.

The research pipeline is deterministic: BLAS is pinned single-thread, all grids are fixed, and no random-number
generator is used. Two independent runs produce a byte-identical result; the core target battery hashes to SHA-256
\texttt{1fb59f55\dots 093d3e92}. The writing phase is gated --- the canonical site cannot be generated until the
research gate is green and the phase is flipped. The seven layer hashes and the release-audit contract are
collected in Appendix~\ref{app:hashes}; the complete \Og\ obstacle ledger in Appendix~\ref{app:irrepro}.

% =====================================================================
\section{Integumentary pathology: thirteen diseases as perturbations}

\answer{Thirteen common skin diseases are reproduced as signed perturbations of the six verified mechanisms.
Opposite entities emerge from opposite drive signs on a single R19 switch, every intervention is the same knob
reversed, and the layer adds no constant while the core-battery hash stays fixed.}

The preceding sections established six verified mechanisms --- the water barrier (T1), wound closure (T2), the
melanin screen (T3), turnover (T4), thermoregulation (T5), and the carcinogenesis kernel. The major skin diseases
follow without new physics: each is one mechanism driven away from its set-point, and each established treatment is
the same control pushed back.

\textbf{Barrier failures (T1).} \emph{Atopic dermatitis} is a chronic erosion of barrier reserve: the intact-basin
depth falls to about 0.55 of healthy and an emollient restores it to about 0.90. \emph{Contact dermatitis} is the
acute mode of the same switch: an external irritant clears the spinodal and the barrier collapses discontinuously
(TEWL jumping about $20\mul$ at insult $\apx$ spinodal), active repair re-crossing to 1.0. \emph{Ichthyosis} sits
on the retention side: a weak desquamation drive holds the corneum near its spinodal, raising stratum-corneum
residence about $6.6\mul$ (critical slowing), reversed by a keratolytic.

\textbf{Turnover off both poles (T4).} \emph{Psoriasis} is turnover past its differentiation spinodal: below
threshold transit is homeostatic ($\apx 28$ days), but above it the exit drive becomes autonomous and
self-accelerating, compressing transit about $20.8\mul$ toward the cited 3--5-day psoriatic window; an
anti-proliferative lowers the exit drive below the spinodal. Psoriasis (acceleration) and ichthyosis (retention)
are the two opposite signs of one switch.

\textbf{Wound non-closure (T2).} A chronic wound is re-epithelialisation held below its own unjamming drive: with
$q^{*}=3.81$ fixed, a free-edge drive of 0.1 stays below the critical 0.15, the sheet stays jammed, and the gap does
not close; raising the drive to 0.25 (debridement, growth factor) carries it over threshold and the wound closes.

\textbf{Pigment off both poles (T3).} \emph{Vitiligo} drives the melanocyte viability switch below its spinodal:
melanin collapses discontinuously to 0.0 and delivered-ultraviolet attenuation returns to 1.0 (no screen),
photoprotection-deficient, repigmentation needing active phototherapy (hysteresis). \emph{Melasma} is the opposite
sign of the same switch: a standing pro-melanogenic drive raises regulated melanin about $3.12\mul$, reversed by
removing the driver. \emph{Albinism} removes the functional screen upstream, so the shared multistage hazard rises
to a relative risk of about $2.58\mul$ versus pigmented skin; an exogenous sunscreen brings it back to about
$1.13\mul$.

\textbf{Thermoregulation off both poles (T5).} \emph{Hypohidrotic ectodermal dysplasia} caps the sweat
appendage's maximum capacity, moving the danger band to lower load (onset $3.1\to1.8$; at a test load of 2.5 the
core measure runs to 4.60 against 2.05 healthy), external cooling capping the load. \emph{Primary hyperhidrosis} is
the opposite sign: the recruitment threshold is lowered (onset 0.2 against 0.3), reversed by threshold-raising
therapy. \emph{Heat stroke} is the capacity limit of the healthy controller: above saturation the regulated
slope (0.77) jumps to the runaway slope (2.00).

\textbf{Oncology.} The single convex multistage rate carries into disease: SCC near-linear ($r=0.999$); melanoma
intermittent-sensitive (burst relative risk up to about $5.67\mul$); a chronic tan protective ($\apx 2.63\mul$);
and the pigment-loss diseases couple in (removing the screen raises the burst relative risk to about $10.59\mul$).
\emph{Actinic keratosis} is the same kernel with fewer hits crossed --- high field prevalence (0.96) above invasive
SCC (0.18) at the same dose.

\begin{table}[h]
\centering\small
\caption{The five-way opposite-sign discriminant: opposite clinical entities from the same R19 switch driven in
opposite directions, no constant changed between poles. All five pairs reproduce \V.}
\begin{tabular}{l l l c}
\toprule
\textbf{Switch} & \textbf{Positive-drive pole} & \textbf{Negative-drive pole} & \textbf{Reproduced}\\
\midrule
Turnover         & psoriasis (accelerated)     & ichthyosis (retention)              & yes\\
Barrier reserve  & intact (full reserve)       & atopic (collapsed reserve)          & yes\\
Melanin          & melasma (overshoot)         & vitiligo (loss)                     & yes\\
Sweat recruitment& hyperhidrosis (excess)      & ectodermal dysplasia (deficit)      & yes\\
Photoprotection  & healthy tan (screened)      & pigment loss (unscreened)           & yes\\
\bottomrule
\end{tabular}
\end{table}

That a single switch generates psoriasis and ichthyosis, melasma and vitiligo, hyperhidrosis and the
ectodermal-dysplasia deficit from opposite signs alone is the strongest internal evidence that these lesions are
dynamics, not bespoke parameters. Every disease mechanism is \V, every clinical mapping rests on a cited anchor
\Lg, and every absolute magnitude is open \Og, each inheriting the obstacle of its parent target. The pathology
layer hashes independently to \texttt{0a4404cc\dots} and does not touch the core-battery hash.

% =====================================================================
\section{The hair-follicle cycle: an emergent oscillator and four alopecias (T6)}

\answer{The hair follicle --- the one autonomously cycling skin structure --- is reproduced as an emergent
relaxation oscillator on the same measured EDAR $\gamma$, with no new fitted constant. The cycle is anagen-dominant
with a plateau-then-collapse waveform; four alopecias follow as signed perturbations of the one cycle.}

Every mechanism so far has been a static R19 switch. But the core battery itself reports a gap --- no autonomous
oscillator anywhere in the integumentary class --- and the structure that demands one is the hair follicle, which
cycles on its own between a long growth phase (anagen), a brief regression (catagen), and a resting phase
(telogen). This is closed with no new organ and no new fitted constant: the \emph{same} measured
$\mathrm{EDAR}\,\gamma=1.3696$ that emerged the skin appendage is placed on the \emph{same} vendored relaxation
element the thermoregulation page already uses, and the cycle is read off. The growth bias is half the appendage
spinodal (0.308), a substrate-scaled regime scale \Fg.

The cycle is anagen-dominant (anagen fraction $\apx 0.59$) with a clear relaxation waveform (the plateau occupies
$\apx 0.87$ of each beat before a fast collapse) \V. The package is deliberately honest about what it does not
claim: the absolute anagen fraction (cited 85--90\%) and the years-long period are \emph{not} fitted \Og, because a
constant drive strong enough to reach 88\% would push the element through its Hopf point and kill the oscillation.
The substrate is therefore graded by dominance and direction, not magnitude.

Four alopecias follow as one oscillator driven off its cycle. \emph{Androgenetic alopecia} is a standing
anti-growth drive: the anagen fraction falls progressively $0.59\to0.54\to0.49$ (miniaturisation), and minoxidil /
anti-androgen lengthens it back to about 0.63 --- the same knob reversed. \emph{Alopecia areata} is a sustained
premature-catagen drive collapsing anagen to about 0.48; removing the drive lets anagen recover to about 0.59
(bistable regrowth hysteresis). The two effluvia share a synchronising cause but differ in timing:
\emph{telogen effluvium} pushes an anagen cohort synchronously into telogen, which sheds exactly one telogen later
--- a shed peak lagging the stressor by 365 steps, precisely $(1-0.59)$ of the cycle --- then self-limits;
\emph{anagen effluvium} is a direct cytotoxic insult to the growing matrix that sheds immediately (lag 0), bypassing
telogen. Same oscillator, opposite timing \V. A three-axis opposite discriminant (anagen short/long; shed
delayed/immediate; loss sustained/self-limited) reproduces on all axes. The layer hashes independently to
\texttt{d910fa5d\dots}; the core hash is unchanged.

% =====================================================================
\section{The sebaceous duct: follicular occlusion as a hysteretic jam (T7)}

\answer{The pilosebaceous duct is reproduced as a hysteretic two-state occlusion jam on the same R19 switch,
running on a newly measured PRDM1 $\gamma$. The duct snaps shut discontinuously at an upper spinodal and reopens
only at a lower one; acne and hidradenitis suppurativa follow as the one jam at two depths.}

The integumentary class is jamming: a barrier or channel driven past a spinodal by an external insult. The
pilosebaceous duct is a textbook member the earlier targets did not cover --- a follicular channel whose net
occlusion drive (androgen-driven sebum, infundibular hyperkeratinisation, and \emph{Cutibacterium acnes}, minus
clearance) can rise until the duct snaps shut into a comedo. It is emerged here as the same R19 jamming switch,
running on a \emph{newly measured} sebaceous master-gene $\gamma$: PRDM1/Blimp1, the committed sebaceous-lineage
master (Horsley et al., 2006), promoted \texttt{to\_measure $\to$ fetched $\to$ vendored} exactly as MITF and EDAR
were. Its $\gamma=1.3432$ is the proximal-promoter nearest-neighbour $\Delta G_{37}$ (SantaLucia, 1998) over the
cached window, computed with the same pipeline validated to reproduce MITF and EDAR byte-for-byte, and reproduces
offline ($\gamma=1.343156$) \Lg. Committing the master gene on the biology before reading $\gamma$ is the
discipline that makes the next result a prediction: PRDM1 carries the lowest $\gamma$ in the atlas, so its functional
spinodal opens earliest --- graded by sign \V.

Run as the R19 switch, the duct does not drift smoothly: it stays patent, then snaps shut discontinuously at the
upper spinodal (a jump of about 2.01 at net occlusion $\apx 1.00$), and reopens not at the same level but only at a
lower spinodal (net occlusion $\apx -1.00$), tracing a hysteresis loop of width about 2.01 \V. The healthy duct
sits patent at net occlusion $-0.30$. The hysteresis \emph{shape} is verified \V; the occlusion set-points are
dimensionless regime scales \Fg; the absolute comedo and lesion counts, Hurley-stage extent, and sebum excretion
rate are open \Og (a per-gland calibration).

\emph{Acne vulgaris} is a standing high occlusion drive (about 1.10) past the upper spinodal, so the duct jams into
a comedo and the \emph{C.\ acnes} amplifier makes it inflammatory; a single agent that only nudges the drive is
insufficient (the duct must be carried below the \emph{lower} spinodal), but a combined comedolytic + sebostatic +
antimicrobial regimen (or isotretinoin) reopens it, while de-inflaming alone leaves a comedonal-but-jammed residue.
\emph{Hidradenitis suppurativa} is the same jam in deeper apocrine-bearing follicles, driven heavier (about 1.50),
where the plug ruptures into the dermis into a chronic scarring sinus-tract sub-state acne never reaches; dropping
the drive to the same sub-acne level that clears acne still leaves the ruptured tract, so a biologic calms activity
but only deroofing/excision resets the scar. A three-way opposite-mode discriminant (reversible superficial /
irreversible deep rupture; inflammatory / comedonal; lighter / heavier plug) reproduces on all axes \V. The layer
hashes independently to \texttt{1e8a557d\dots}; the core hash is unchanged.

% =====================================================================
\section{Cell adhesion as a hysteretic binding jam: pemphigus and pemphigoid (T8)}

\answer{Junctional adhesion is reproduced as a hysteretic two-state binding jam on the already-measured KRT14
$\gamma$, with no new $\gamma$ fetched and none fitted. Pemphigus vulgaris and bullous pemphigoid differ only by
which compartment the antibody hits, flipping the cleavage plane, the Nikolsky sign, and the blister tension.}

Desmosomes and hemidesmosomes anchor the keratinocyte's keratin network, so junctional adhesion is an intrinsic
property of the keratinocyte --- a new \emph{target} on an existing organ (the hair-cycle pattern), not a new
measured organ. It runs on the already-measured $\mathrm{KRT14}\,\gamma=1.4894$ with no new $\gamma$ fetched and
none fitted. A bond is read the natural way round: \emph{adherent} $=$ jammed ON, \emph{blistered} $=$ unjammed OFF.
Run as the switch, the bond detaches discontinuously at the lower spinodal (a jump of about 2.11 at net adhesion
$\apx -1.005$) and re-adheres only at a higher spinodal (net adhesion $\apx +1.005$), tracing a hysteresis loop of
width about 2.01 \V; the healthy bond sits above the upper spinodal at net adhesion $+1.40$. Adhesion runs in two
coupled compartments on the \emph{same} $\gamma$ and spinodal: cell-cell (desmosomal, DSG3) and cell-matrix
(hemidesmosomal, BP180/COL17A1), and an antibody of the same magnitude detaches whichever it is specific for.

\emph{Pemphigus vulgaris}: an anti-DSG3 autoantibody drives the cell-cell compartment below the spinodal $\to$ an
intraepidermal/suprabasal split with the cell-matrix bond intact (basal ``tombstone'' row). Because the failing
bond is the lateral one, a tangential shear propagates, so the Nikolsky sign is \emph{positive} (derived, not
asserted) and the thin roof is flaccid; a partial titre reduction does not re-adhere, clearance (rituximab /
immunosuppression) does. \emph{Bullous pemphigoid}: an anti-BP180 antibody of the \emph{same magnitude} instead
drives the cell-matrix compartment below the spinodal $\to$ a subepidermal split with cell-cell bonds intact (the
epidermis lifts whole). The failing bond is basal, not lateral, so shear does not propagate: the Nikolsky sign is
\emph{negative} (derived) and the full-thickness roof is tense; corticosteroid / immunosuppression re-adheres it.

\begin{table}[h]
\centering\small
\caption{The three-axis opposite-property discriminant for autoimmune blistering: from the targeted
\emph{compartment} alone, at the same antibody magnitude on one reused KRT14 $\gamma$. All three pairs reproduce \V.}
\begin{tabular}{l l l c}
\toprule
\textbf{Axis} & \textbf{Pemphigus vulgaris} & \textbf{Bullous pemphigoid} & \textbf{Reproduced}\\
\midrule
Cleavage plane & intraepidermal / suprabasal & subepidermal / junctional & yes\\
Nikolsky sign  & positive (derived)          & negative (derived)        & yes\\
Blister tension& flaccid                     & tense                     & yes\\
\bottomrule
\end{tabular}
\end{table}

The adhesion reserve and antibody titres are dimensionless regime scales \Fg; the clinical mappings are cited
anchors \Lg; absolute blister counts, titre (IU/mL), cleavage depth (\um), and body-surface area are open \Og (a
per-junction calibration). Boundary: congenital epidermolysis bullosa (a defective adhesion \emph{gene}
KRT14/COL17A1/LAMB3) is gene-keyed and belongs to the gene-lesion registry, not this dynamical layer. The layer
hashes independently to \texttt{55dce8c2\dots}; the core hash is unchanged.

% =====================================================================
\section{Neurovascular reactivity as a hysteretic vasomotor jam: rosacea and Raynaud (T9)}

\answer{Cutaneous vasomotor tone is reproduced as a hysteretic two-lock reactivity jam on the already-measured EDAR
$\gamma$, with the dermal-perfusion magnitude left an inherited circulatory seam. Rosacea (a vasodilator drive that
locks the vessel dilated) and Raynaud phenomenon (a vasoconstrictor drive giving a reversible attack) differ only by
drive sign.}

Every disease so far lived on a barrier, a bond, a duct, or a cycle. Rosacea lives on the cutaneous vasculature's
reactivity, and the core battery had no vasomotor target --- which is exactly why it was held back. The target is
added without a new organ and without a fitted constant: the already-measured $\mathrm{EDAR}\,\gamma=1.3696$ --- the
thermoregulation-interface organ --- is placed on the same R19 switch. The justification is anatomical: the
cutaneous thermoregulatory interface carries two autonomic effector arms, the sudomotor arm (sweat, the T5 target)
and the vasomotor arm (skin blood flow); rosacea dysregulates the vasomotor arm, the \emph{vascular mirror of
hyperhidrosis} on the sudomotor arm, so it is intrinsic to this organ. A vessel is read the natural way round:
\emph{dilated} $=$ jammed ON, \emph{constricted} $=$ OFF. The dermal-perfusion \emph{magnitude} is not re-derived:
it stays an inherited circulatory citation; only the reactivity dynamics is added.

Run as the switch, the vessel locks dilated discontinuously once the net dilator drive exceeds the upper spinodal
(a jump of about 2.03 at net drive $\apx +1.005$) and locks constricted discontinuously once it falls below the
lower spinodal (at net drive $\apx -1.005$), tracing a hysteresis loop of width about 2.01 \V; between the locks it
is reversibly responsive (healthy net drive $-0.30$). Reversibility is a uniform consequence of \emph{whether a
drive crosses its lock}: an early flush comes and goes, but once a sustained drive carries the vessel past the upper
lock the dilation is fixed.

\emph{Rosacea} is a standing vasodilator reactivity drive (a lowered flush threshold under heat / alcohol /
ultraviolet / spice) that carries the net dilator drive past the upper spinodal, so the vessel locks dilated into
persistent erythema and fixed telangiectasia (dilation depth 1.00); the LL-37 / Demodex inflammatory amplifier on
the same dilated background gives the papulopustular subtype. Anti-inflammatory therapy clears the papules but
leaves the vascular background; an $\alpha$-agonist (brimonidine) blanches transiently but, because a partial
constrictor leaves the net drive above the lower lock, does not reset the fixed vessels (hysteresis); only laser /
intense-pulsed-light physically resets the telangiectasia --- a structural reset exactly like deroofing for a
hidradenitis tract. \emph{Raynaud phenomenon} is the opposite sign: a cold / stress vasoconstrictor drive carries
the net drive negative into the constricted/ischemic basin (constriction depth 1.00) producing the white-blue-red
digital attack; \emph{primary} Raynaud crosses no lock, so it reverses on rewarming and a vasodilator /
calcium-channel blocker (nifedipine) aborts it. The \emph{fixed} digital-ischemia / ulcer state of \emph{secondary}
Raynaud (connective-tissue disease) is a downstream structural change --- an immune/rheumatology sibling-package
seam, named not modelled.

\begin{table}[h]
\centering\footnotesize
\caption{The three-axis opposite-property discriminant for the vasomotor diseases: from the drive sign and the lock
rule alone, on one reused EDAR $\gamma$. All three pairs reproduce \V.}
\begin{tabular}{@{}p{2.9cm} p{4.5cm} p{4.9cm} c@{}}
\toprule
\textbf{Axis} & \textbf{Rosacea} & \textbf{Raynaud phenomenon} & \textbf{Repr.}\\
\midrule
Direction              & vasodilation (above upper lock)   & vasoconstriction (into constricted basin) & yes\\
Reversibility (lock)   & fixed telangiectasia (crosses lock)& reversible vasospasm (short of lock)      & yes\\
Vascular vs.\ inflammatory & erythematotelangiectatic       & papulopustular (LL-37/Demodex amplifier)  & yes\\
\bottomrule
\end{tabular}
\end{table}

The resting tone, reactivity gains, and constrictor drives are dimensionless regime scales \Fg; the clinical
mappings are cited anchors \Lg; absolute erythema index, vessel density, flush magnitude, digital temperature,
attack frequency, and body-surface area are open \Og (a per-vessel calibration), with the perfusion magnitude an
inherited circulatory seam. The layer hashes independently to \texttt{53a99f52\dots}; the core hash is unchanged.

% =====================================================================
\section{Cross-package seams: what the package imports, exposes, and hands off}

\answer{For the VP organs to run as one body, each package imports a sibling's result across a named seam rather
than re-deriving it. This is the integumentary package's labelled seam record: it adds no new mechanism, exposes the
one internal cross-target coupling (pigment-loss removes the melanin screen and raises the oncology hazard), and
declares what it inherits and hands off.}

The whole VP programme is a family of organ packages that must eventually run as one body without any package
secretly re-deriving another's result. The discipline that makes that safe is the single source of truth (SSOT):
each quantity is owned by exactly one package and imported elsewhere across a named seam. This section is the
integumentary package's complete, machine-readable seam record (also emitted as
\texttt{reports/seam\_manifest.json}), in three honestly distinct classes; it adds no new mechanism and no new
constant.

\textbf{Inherited-in (3).} Read-only inputs vendored from a sibling and never re-derived: the circulatory
dermal-perfusion \emph{magnitude} (cited, consumed by T9); the DNA organ identity, emergence order, and measured
$\gamma$ \V; and the R19 substrate primitive.

\textbf{Internal-live (1).} The one genuinely internal cross-target coupling: a melanocyte-target (T3) lesion that
removes the melanin screen raises the shared oncology-kernel hazard. The numbers are \emph{re-exported verbatim}
from the verified pathology layer, so the seam output provably \emph{is} the internal link surfaced, not a parallel
implementation. Removing the screen with an albinism-type lesion raises the SCC cumulative-hazard relative risk to
about $2.58\mul$ (incidence relative risk about $2.24\mul$) and the melanoma \emph{burst} relative risk to about
$10.59\mul$; restoring an exogenous screen (sunscreen) brings the SCC hazard back to about $1.13\mul$ --- the screen
is the causal lever. (That the lesion's delivered-ultraviolet attenuation is exactly 1.00 --- no screen at all ---
is the mechanistic reason.) The coupling shape is \V, the epidemiology cited \Lg, the absolute relative risk open
\Og. See \S\ref{sec:anchors} for the honesty flag: the $10.59\mul$ melanoma-burst figure is a property of the shared
oncology \emph{kernel}, not a claim that albinos have $\apx 10.6\mul$ the melanoma incidence (OCA melanoma is
clinically rare because the melanocytic substrate is removed --- consistent, on a different axis).

\textbf{Declared-out (6).} Contracts to sibling packages, flagged honestly as \emph{declared, not yet live
wiring} --- the integration harness does not exist in this self-contained package. They are: gene-lesion $\to$
\texttt{disease\_wp} (single-gene genodermatoses); urticaria/angioedema (immune-hematologic); lichen planus /
inflammatory dermatoses (immune-effector); the \emph{fixed} digital ischemia of secondary Raynaud (rheumatology);
the rosacea inflammatory chemistry beyond the LL-37/Demodex flag (immune-effector); and out-of-class cutaneous
infections (impetigo, cellulitis, dermatophytosis, herpes, warts), which are pathogen-driven, not an R19-dynamics
failure. Flagging them as declared rather than silently modelling them is the same honesty rule as grading an
absolute magnitude \Og: the boundary is named, not hidden.

% =====================================================================
\appendix
\section*{\centering\rule{0.4\linewidth}{0.4pt}\ \ Appendices\ \ \rule{0.4\linewidth}{0.4pt}}
\addcontentsline{toc}{section}{Appendices}

% ---------------------------------------------------------------
\section{The frozen determinism contract}\label{app:hashes}

The v1.0.0 release is label-only over v0.9.0's computation: the version string and build date were bumped and the
HTML re-stamped, while the seven hashes below are byte-identical to v0.9.0. The release adds exactly one additive,
read-only artifact (\texttt{repro/run\_release\_audit.py}) and the governance documents. The evidence of record is
these seven hashes plus the green gates; \texttt{python repro/run\_release\_audit.py} re-runs each canonical runner
in its own subprocess and asserts each hash matches the pinned value. The verdict on this release is $7/7$ hashes
match, all gates green, \textbf{release-ready}.

\begin{table}[h]
\centering\small
\caption{Coverage certified by the release-readiness audit (\texttt{reports/release\_audit.json},
\texttt{all\_ok = true}).}
\begin{tabular}{l l}
\toprule
\textbf{Check} & \textbf{Value}\\
\midrule
In-lane program (jamming class) & complete\\
Targets verified & 10 \;(T1--T9 + oncology kernel)\\
Diseases covered + verified & 23 \;($13+4+2+2+2$)\\
Seam manifest & inherited-in 3 / internal-live 1 / declared-out 6\\
Canonical doc sections (HTML, C4) & 14 + hub\\
Frozen determinism hashes & $7/7$ match\\
All layer gates green & yes\\
No-tuning discipline & held (no constant fitted; \texttt{\_to\_measure} empty)\\
\bottomrule
\end{tabular}
\end{table}

\begingroup
\footnotesize\ttfamily\setlength{\parskip}{0.15em}
\textbf{\rmfamily Layer --- runner --- SHA-256 (two independent runs identical):}\par
\smallskip
core T1--T5 + oncology --- run\_all.py\par
\hangindent=2em 1fb59f556e01882916c0f4b4e72aa76ad47ef348bffb81fe4a7f3ce9093d3e92\par
pathology (13 + 5-way) --- run\_pathology.py\par
\hangindent=2em 0a4404ccde6559a7621249de90cd889f907619db48a4104329f3af2f74c51cd8\par
hair-cycle (T6 + 4) --- run\_cycle.py\par
\hangindent=2em d910fa5d2854462236a5ef56b831490f1aca82ba37050b3f8bb751198b356822\par
sebaceous (T7 + 2) --- run\_seb.py\par
\hangindent=2em 1e8a557d9a8d7b05823259bb2fcac31e372b96272e47ef1f2f5d89b1b0246a84\par
adhesion (T8 + 2) --- run\_adhesion.py\par
\hangindent=2em 55dce8c267ea4c76d8d5b967535a4ee0876c4b67f605e6ae7d65bf38516f28ad\par
vasomotor (T9 + 2) --- run\_vasomotor.py\par
\hangindent=2em 53a99f522ad684a11bcc6d4d33c5e8123d1f4dff5a75836c0bbced6736a62003\par
seam manifest --- run\_seam.py\par
\hangindent=2em 52b49a95a9add070a05a02848b1a4cef0589f4a36791167232fda3fc1c1d58f7\par
\endgroup

% ---------------------------------------------------------------
\section{Irreproducibility ledger (the \texorpdfstring{\Og}{[O]} obstacle list)}\label{app:irrepro}

Every open quantity is listed with the specific external input that would close it (VP-SPEC C3). The shapes,
controls, orderings, and dichotomies in the body stand independently of these open absolute magnitudes.

\begin{longtable}{p{4.6cm} c p{8.6cm}}
\toprule
\textbf{Item} & \textbf{Grade} & \textbf{Obstacle (why not reproducible in-package)}\\
\midrule
\endhead
to-measure master-gene $\gamma$ & \Og$\to$\V & none deferred (\texttt{\_to\_measure: []}); any future master is fetched via the DNA atlas as a measured input, not fitted\\
absolute TEWL (\tewlunit) & \Og & T1 threshold + discontinuous-collapse shape \V; absolute flux needs stratum-corneum lipid permeability and the trans-barrier water-activity gradient\\
absolute wound closure rate (\ensuremath{\mathrm{\mu m\,h^{-1}}}) & \Og & T2 unjam$\to$migrate$\to$rejam sequence + chronic-wound threshold \V; $q^{*}=3.81$ cited \Lg; absolute rate needs a measured single-cell migration speed\\
absolute minimal erythema dose / melanin OD & \Og & T3 concave-plateau response, partial attenuation, feedback-off control \V; absolute MED needs the melanin extinction coefficient\\
absolute basal cycle time (per cell) & \Og & T4 conveyor-sum structure \V and rate set by one cited anchor (SC transit $\apx 14$ d) \Lg; per-cell cycle needs a per-cell calibration\\
absolute set-point (37\dg C) \& sweat rate & \Og & T5 thresholded onset, slope drop, runaway-above-capacity \V; absolute set-point/rate need per-gland output and body heat capacity\\
absolute cancer incidence \& RR magnitude & \Og & oncology dichotomy directions \V on cited anchors \Lg (Gandini SRR 1.61; Armitage--Doll $K\apx5$); absolute incidence/RR need a population baseline rate and absolute dose calibration\\
absolute anagen fraction (cited 85--90\%) \& period (years) & \Og & anagen dominance and the relaxation waveform are dimensionless \V; the absolute fraction cannot be tuned without crossing the Hopf point and killing the limit cycle (per-follicle calibration)\\
absolute comedo/lesion counts; Hurley extent; sebum rate & \Og & T7 occlusion-jam direction, discontinuity, hysteresis, reopening \V; set-points are regime scales \Fg; absolute counts need a per-gland calibration\\
absolute blister counts; titre (IU/mL); cleavage depth (\um); BSA & \Og & T8 de-adhesion direction, plane selection, derived Nikolsky sign, re-adhesion \V; titres are regime scales \Fg; absolutes need a per-junction calibration\\
absolute erythema index; vessel density; digital temperature; BSA & \Og & T9 lock rule, both discontinuous locks, reversible middle \V; drives are regime scales \Fg; absolutes need a per-vessel calibration; the perfusion magnitude is an inherited circulatory seam\\
absolute organ size / mass & \Og & dwell $\propto \gamma^{1.5}$ fixes relative size/order \Fg; absolute scale needs an external calibration\\
\bottomrule
\end{longtable}

% ---------------------------------------------------------------
\section{Anchor verification (v1.0.0, against current literature)}\label{sec:anchors}

A v1.0.0 pass over the cited \Lg\ anchors against current authoritative literature. It changes no computation:
every determinism hash is byte-frozen; the \V\ shapes were never fitted to these numbers; the \Og\ magnitudes
remain open with their stated obstacle. Its job is to make each citation explicit and to flag, with the framework's
``name it, don't hide it'' discipline, the one place where a model quantity must not be read as a literal incidence.

\begin{itemize}[leftmargin=1.4em,itemsep=2pt]
\item \textbf{Epidermal turnover (T4).} Current literature: stratum-corneum transit $\apx 14$--20 d;
whole-epidermis turnover $\apx 28$--48 d (classic Weinstein-class), extending to $\apx 56$ d in newer density-based
estimates; the widely-quoted ``28-day renewal'' is a lower-bound simplification. The modelled $\apx 28$ d total
therefore sits at the \emph{conservative lower edge} of the literature --- stated, not tuned; the \V\ dwell-cascade
structure is unchanged, the absolute total stays \Og. The CHARTER \texttt{TO-ANCHOR} flag is cleared.
\item \textbf{UV dichotomy.} Confirmed: melanoma $\leftarrow$ intermittent/recreational + sunburn (Gandini summary
relative risk $\apx 1.6$), with a chronic/occupational \emph{inverse} association (the tan paradox); SCC
$\leftarrow$ cumulative ultraviolet, near-linear (occupational pooled odds ratio $\apx 1.77$). Matches the model's
\V\ shapes. Refinement recorded: BCC is more intermittent-like than SCC; a future SCC/BCC split could place BCC
nearer the melanoma pole.
\item \textbf{Pigment-loss $\to$ oncology (the internal-live seam).} OCA $\to$ SCC confirmed (SCC dominant,
$\apx 75$--88\% of albino skin cancers; up to $\apx 1000\mul$ risk in sub-Saharan-African albinos; sunscreen the
preventive/causal lever). \emph{Honesty flag}: the seam's ``melanoma burst relative risk $\apx 10.59\mul$'' is a
property of the shared oncology kernel (screen removed $\to$ the multistage rate's burst-sensitivity rises), not a
claim that albinos have $\apx 10.6\mul$ the melanoma incidence --- OCA melanoma is clinically rare because OCA
removes the melanocytic substrate. The two are consistent (different axes), and the package already grades the burst
relative risk \Og, never as a cited incidence. No computation changed.
\end{itemize}

Net effect on the release: no hash changed, no constant added, no shape re-fitted. The \Lg\ anchors are now
explicit and current, the \texttt{TO-ANCHOR} flag is cleared, and the single place where a model number could be
mis-read epidemiologically is named and bounded.

% ---------------------------------------------------------------
\section{How to reproduce}\label{app:reproduce}

The reproducible engine is bundled in the accompanying archive under \texttt{repro/}. With Python 3 and NumPy
(BLAS is pinned single-thread inside the engine), from the package root:

\begingroup\small\ttfamily
\setlength{\parskip}{0.1em}
python repro/run\_all.py \hfill \rmfamily\itshape core battery (T1--T5 + oncology) --- sha 1fb59f55\dots\par
python repro/run\_pathology.py \hfill \rmfamily\itshape 13 diseases + 5-way discriminant --- sha 0a4404cc\dots\par
python repro/run\_cycle.py \hfill \rmfamily\itshape hair-cycle oscillator + 4 alopecias --- sha d910fa5d\dots\par
python repro/run\_seb.py \hfill \rmfamily\itshape sebaceous-duct jam + 2 diseases --- sha 1e8a557d\dots\par
python repro/run\_adhesion.py \hfill \rmfamily\itshape cell-adhesion jam + 2 diseases --- sha 55dce8c2\dots\par
python repro/run\_vasomotor.py \hfill \rmfamily\itshape vasomotor jam + 2 diseases --- sha 53a99f52\dots\par
python repro/run\_seam.py \hfill \rmfamily\itshape cross-package seam manifest --- sha 52b49a95\dots\par
python repro/run\_release\_audit.py \hfill \rmfamily\itshape re-runs all of the above; asserts $7/7$ hashes\par
\endgroup

\medskip
\begin{sloppypar}
Each runner prints its $2\times$ SHA-256 and writes a machine-readable record under \texttt{reports/}. The
additive layers (\texttt{run\_pathology}, \texttt{run\_cycle}, \texttt{run\_seb}, \texttt{run\_adhesion},
\texttt{run\_vasomotor}, \texttt{run\_seam}) each carry their own hash and never touch the core hash. The canonical
HTML is rebuilt by \texttt{python tools/build\_docs.py}, which refuses to run unless the research gate is green and
\texttt{PHASE = writing}; every displayed number is pulled live from the engine.
\end{sloppypar}

% =====================================================================
\section*{Selected cited anchors}
\addcontentsline{toc}{section}{Selected cited anchors}
\small

The \Lg-graded external anchors used in the body. All \V\ shapes are derived from the substrate and were never
fitted to these numbers.

\begin{enumerate}[leftmargin=1.6em,itemsep=1.5pt]
\item D.~Bi, J.~H.~Lopez, J.~M.~Schwarz, M.~L.~Manning. \textit{A density-independent rigidity transition in
biological tissues.} Nature Physics \textbf{11}, 1074--1079 (2015); and \textit{Motility-driven glass and jamming
transitions in biological tissues.} Phys.\ Rev.\ X \textbf{6}, 021011 (2016). --- the vertex-model shape-index
threshold $q^{*}\apx3.81$ (T2).
\item S.~Gandini et al. \textit{Meta-analysis of risk factors for cutaneous melanoma: II.\ Sun exposure.}
European Journal of Cancer \textbf{41}, 45--60 (2005). --- intermittent summary relative risk 1.61; chronic
exposure inversely associated; sunburn risk (T3 oncology, \S7, \S9).
\item P.~Armitage, R.~Doll. \textit{The age distribution of cancer and a multi-stage theory of carcinogenesis.}
British Journal of Cancer \textbf{8}, 1--12 (1954). --- the multistage exponent $K\apx5$ (oncology kernel).
\item J.~SantaLucia~Jr. \textit{A unified view of polymer, dumbbell, and oligonucleotide DNA nearest-neighbor
thermodynamics.} PNAS \textbf{95}, 1460--1465 (1998). --- the nearest-neighbour $\Delta G_{37}$ used to compute
master-gene $\gamma$ from cached promoters (PRDM1, T7).
\item V.~Horsley et al. \textit{Blimp1 defines a progenitor population that governs cellular input to the sebaceous
gland.} Cell \textbf{126}, 597--609 (2006). --- PRDM1/Blimp1 as the committed sebaceous-lineage master (T7).
\item G.~D.~Weinstein, J.~L.~McCullough, P.~Ross. \textit{Cell proliferation in human epidermis.} J.\ Invest.\
Dermatol. (classic radiolabel turnover; stratum-corneum transit and whole-epidermis renewal windows) (T4,
\S\ref{sec:anchors}).
\end{enumerate}

\vfill
\begin{center}\footnotesize\itshape
Single source of truth: the canonical artifact is the per-title HTML at jamming-physics.org/integumentary.\\
This whitepaper and the bundled engine are reproducible to the seven SHA-256 hashes above.\\
Young Jae Lee --- ORCID 0009-0002-7535-8245 --- VP Theory --- CC BY 4.0 --- v1.0.0, 2026-06-19.
\end{center}

\end{document}
